\section{Conclusion}
\label{sec:conclusion}
%% ===================================================================

Lux is not a database. Lux is a trust kernel.

This paper has formalized the trust kernel abstraction: the minimum on-chain
substrate that makes sovereign applications composable, auditable, and
provider-independent. The four chains---K (key policy), T (threshold governance),
M (MPC signing), I (decentralized identity)---provide exactly the trust services
that local encrypted vaults cannot provide on their own: global identity,
cross-party coordination, and settlement finality.

The trust kernel stores no application data. It stores commitments to
decisions about data. Every decision produces a receipt. Receipts are
permanent, tamper-evident, and verifiable by any party. This architecture
enables capsules to switch providers, delegate capabilities to other capsules,
and participate in threshold governance---all without revealing plaintext to
the chain or to each other.

The post-quantum migration path ensures that trust decisions made today remain
verifiable under future quantum adversaries. The mechanized Lean~4 proofs
provide formal guarantees for the critical protocol properties: consensus
safety, signature unforgeability, threshold composition security, and trust
delegation integrity.

The trust kernel is the answer to a question that the blockchain industry has
been asking incorrectly. The question is not ``how do we put data on-chain.''
The question is ``what is the minimum we must put on-chain so that everything
else can stay off-chain and still be trusted.'' Four things: identity, key
policy, threshold governance, and receipts.

%% ===================================================================
%% References
%% ===================================================================

\begin{thebibliography}{99}

\bibitem{lux-capsule-architecture}
Z.~Kelling.
\newblock Lux Capsule Architecture: Sovereign Encrypted Vaults with Trust Kernel Binding.
\newblock Lux Industries, 2026.

\bibitem{lux-mchain-mpc}
Z.~Kelling and V.~Seesahai.
\newblock M-Chain: Decentralized Multi-Party Computation Custody with Quantum-Safe Threshold Signatures.
\newblock Lux Industries, 2023.

\bibitem{lux-id-did-specification}
Z.~Kelling.
\newblock Lux ID: Decentralized Identity and Universal Access Management.
\newblock Lux Industries, 2020.

\bibitem{lux-pq-crypto-suite}
Z.~Kelling.
\newblock Post-Quantum Cryptographic Suite for EVM: ML-KEM, ML-DSA, and SLH-DSA as Native Precompiles.
\newblock Lux Industries, 2024.

\bibitem{lux-tfhe}
Z.~Kelling.
\newblock Torus Threshold Fully Homomorphic Encryption: Boolean Circuits on Encrypted Data with Distributed Decryption.
\newblock Lux Industries, 2025.

\bibitem{lux-warp-messaging}
Z.~Kelling.
\newblock Lux Warp Messaging: Cross-Chain Communication Protocol.
\newblock Lux Industries, 2023.

\bibitem{ucan-fission}
B.~Frazee, P.~Willison, and B.~Holmgren.
\newblock UCAN: User Controlled Authorization Networks.
\newblock Fission Codes, 2021.

\bibitem{ceramic-anchors}
J.~Clark, M.~Zargham, and D.~Rees.
\newblock Ceramic: A Decentralized Data Network for Web3.
\newblock 3Box Labs, 2022.

\bibitem{keybase}
M.~Krohn, M.~Mazieres, and C.~Coyne.
\newblock Keybase: Public Key Crypto for Everyone, Publicly Auditable Proofs of Identity.
\newblock Keybase, Inc., 2014.

\bibitem{signal-protocol}
T.~Perrin and M.~Marlinspike.
\newblock The X3DH Key Agreement Protocol.
\newblock Signal Foundation, 2016.

\bibitem{frost}
C.~Komlo and I.~Goldberg.
\newblock FROST: Flexible Round-Optimized Schnorr Threshold Signatures.
\newblock In \emph{Selected Areas in Cryptography (SAC)}, 2020.

\bibitem{cggmp21}
R.~Canetti, R.~Gennaro, S.~Goldfeder, N.~Makriyannis, and U.~Peled.
\newblock UC Non-Interactive, Proactive, Threshold ECDSA with Identifiable Aborts.
\newblock In \emph{ACM CCS}, 2020.

\bibitem{mlkem}
NIST.
\newblock FIPS 203: Module-Lattice-Based Key-Encapsulation Mechanism Standard.
\newblock National Institute of Standards and Technology, 2024.

\bibitem{mldsa}
NIST.
\newblock FIPS 204: Module-Lattice-Based Digital Signature Standard.
\newblock National Institute of Standards and Technology, 2024.

\end{thebibliography}

